We compare four training objectives for a small CNN under symmetric label noise on the scikit-learn 8$\times$8 digits: cross-entropy (CE), label smoothing (LS), mixup, and a one-network small-loss selection rule in the spirit of co-teaching, at 0, 20, 40 and 60\% noise with 5 seeds each, all reimplemented with identical architecture, schedule and data. Besides clean test accuracy, we measure memorisation on the corrupted training examples (the fraction fitted to the wrong label minus the fraction predicted as the true label). CE memorises nearly all corrupted labels and drops from 0.991 to 0.709 at 40\% noise. LS does not help (0.668 at 40\%). Mixup with $\alpha{=}1$ helps modestly (0.762) yet still memorises 89\% of corrupted labels. Small-loss selection, given the true noise rate, is best at 20--60\% noise (0.896 at 40\%, 0.771 at 60\%). Sensitivity sweeps show that under-estimating the noise rate erodes its gain, that over-estimating it does not, and that larger mixup $\alpha$ keeps improving. Test accuracy logged at earlier epochs shows LS above CE and mixup close to small-loss before memorisation sets in, so these rankings are specific to final-epoch evaluation. The study is small (one tiny dataset, untuned baselines, oracle forget rate, final-epoch evaluation), so the ranking should be read as a controlled small-scale observation, not as a general result.
